\section{Range minimum query chaining}
\label{sec:chaining-rmq-idea}

Chaining needs to choose a predecessor for each anchor \(i\). A long search
window may contain many candidates, so calculating every complete chain score
\(f_j+E(j,i)\) can be expensive. Range minimum query (RMQ) chaining first
selects one candidate using a simpler value, then calculates its complete
score. Following the linear scoring idea introduced in
Equation~\ref{eq:global-score-definition-chaining-rmq}, minimap2 stores the
\emph{ranking value}
\begin{equation}
  h(j)=f_j+\frac{\lambda_g}{2}(r_j+q_j).
\label{eq:chaining-rmq-h}
\end{equation}
The important distinction is between selecting a candidate and accepting its
link:
\begin{itemize}
  \item \textbf{Select by \(h(j)\).} Once \(f_j\) is known, \(h(j)\) depends
    only on anchor \(j\). A tree can reuse it for later targets. It stores
    priority \(\pi(j)=-h(j)\), so the smallest priority gives the largest
    ranking value.
  \item \textbf{Check the complete score.} The value \(h(j)\) does not include
    the overlap reward or the full gap penalty for the pair \((j,i)\).
    Its largest value need not give the largest \(f_j+E(j,i)\).
  \item \textbf{Also check nearby anchors.} Minimap2 can scan a smaller
    interval using the complete score. This gives nearby candidates another
    chance even when the tree selected a different anchor.
\end{itemize}
The minigraph paper describes this combination of local scoring and a wider
RMQ search~\cite{minimap2_rmq_scoring_minigraph}.

\subsection{Minimap2 RMQ}
\label{sec:chaining-minimap2-rmq}

The search uses the following limits:
\begin{itemize}
  \item \(D\) is the maximum reference distance for the primary tree and
    sets the width of its query interval. Minimap2 raises the requested
    distance to at least the bandwidth \(B\).
  \item \(D_{\mathrm{in}}\) is the local distance limit, kept between
    \(0\) and \(D\). Zero disables the local scan.
  \item \(B\) is the maximum allowed diagonal difference
    \(\delta(j,i)=|d_r(j,i)-d_q(j,i)|\).
  \item \(C\) limits the number of retained anchors in each tree.
    \(S\) is the marked-anchor skip limit used during the local scan.
\end{itemize}

The primary tree \(T\) keeps earlier anchors on the same reference and strand
with \(0<r_i-r_j\leq D\). Anchors at the current reference coordinate are
inserted only after that whole coordinate group has been scored. Old anchors
are removed when they leave the distance window or when the tree exceeds
\(C\). If \(D_{\mathrm{in}}>0\), a second tree \(T_{\mathrm{in}}\) follows
the same rules with the smaller reference distance \(D_{\mathrm{in}}\).

Both trees are balanced binary search trees (AVL trees). Their search key
and stored priority are
\[
  K_j=(q_j,j),
  \qquad
  \pi(j)=-h(j).
\]
Keys are compared first by query coordinate, then by anchor index.
The index \(j\) keeps equal query coordinates as separate nodes. Each node
also caches the anchor with the smallest priority in its subtree.  The query uses the closed
interval between
\begin{equation}
  K_i^{\mathrm{lo}}=(q_i-D,\mathrm{INT32\_MAX}),
  \qquad
  K_i^{\mathrm{hi}}=(q_i,0).
\label{eq:chaining-rmq-key-range}
\end{equation}
For ordinary positive indices, these composite keys admit
\(q_i-D<q_j<q_i\). The upper key also admits index zero at \(q_j=q_i\).
The intended predecessor relation requires \(q_j<q_i\).
The tree query does not define a stable index tie rule when two anchors have
the same \(h(j)\).

Let \(j^\star\) be the tree's selected candidate. Minimap2 calculates
\(c_{j^\star}=f_{j^\star}+E(j^\star,i)\) and accepts it only when
\(\delta(j^\star,i)\leq B\) and \(c_{j^\star}>f_i\). It calls the link
\emph{exact} when
\begin{equation}
  \delta(j^\star,i)=0
  \quad\land\quad
  g(j^\star,i)\leq\ell_{j^\star}.
  \label{eq:chaining-rmq-exact}
\end{equation}
This means that the two coordinates advance equally and by no more than the
predecessor's seed span. The link has no gap penalty. Here, \emph{exact}
names a condition for skipping further search. It does not prove that this
predecessor gives the largest complete score or the correct alignment.

If the selected link is not exact, the inner tree is nonempty, and
\(q_i>0\), minimap2 calls \textsc{ScanLocalInterval}. This also happens when
the selected link failed the bandwidth or score check:
\begin{itemize}
  \item \textbf{Local candidates.} The scan visits retained anchors with
    \[
      0<r_i-r_j\leq D_{\mathrm{in}},\qquad
      q_i-D_{\mathrm{in}}\leq q_j\leq q_i-1.
    \]
    It reads them by decreasing \((q_j,j)\), so larger query coordinates
    come first. Equal query coordinates use the larger index first.
  \item \textbf{Score updates.} It starts with the current \((f_i,p_i)\),
    including any accepted tree result. For every visited \(j\), it
    calculates \(c_j=f_j+E(j,i)\). It stores \((c_j,j)\) only if
    \(\delta(j,i)\leq B\) and \(c_j>f_i\). An equal score keeps the
    predecessor already saved.
  \item \textbf{Early stopping.} The scan uses the B1 marks and skip
    counter. A marked candidate that passes the bandwidth test but does not
    improve the score increases the counter. An improvement decreases a
    positive counter by one. The scan stops when the counter exceeds \(S\)
    or the local interval ends.
\end{itemize}
Thus the local scan may test several candidates, but need not visit the
whole interval. It does not ask the primary tree for its second choice.
Figure~\ref{fig:chaining-minimap2-rmq} shows the search regions and the checks.

\begin{figure}[H]
\centering
\begin{tikzpicture}[x=1cm,y=1cm,font=\scriptsize,>=Latex]
  \node[anchor=west,font=\bfseries] at (.05,6.35)
    {(a) Maintain the active set};
  \node[anchor=west,font=\bfseries] at (4.85,6.35)
    {(b) Query cached values};
  \node[anchor=west,font=\bfseries] at (10.05,6.35)
    {(c) Local scan and checks};

  % Active reference window.
  \draw[fill=chainGray!10,draw=chainGray!80!black]
    (.15,4.85) rectangle (1.00,5.55);
  \node at (.575,5.20) {\(j_0\)};
  \draw[very thick,draw=chainOrange!95!black]
    (.10,4.78)--(1.05,5.62);
  \node[anchor=north,text=chainGray!90!black] at (.575,4.72) {expired};

  \foreach \x/\lab in {1.25/\(j_1\),2.20/\(j_2\)}
    {
      \draw[fill=chainBlue!13,draw=chainBlue!90!black]
        (\x,4.85) rectangle (\x+.85,5.55);
      \node at (\x+.425,5.20) {\lab};
    }
  \draw[fill=chainGreen!14,draw=chainGreen!80!black,very thick]
    (4.05,4.85) rectangle (4.70,5.55);
  \node at (4.375,5.20) {\(i\)};
  \draw[fill=chainOrange!14,draw=chainOrange!95!black,dashed]
    (4.05,3.75) rectangle (4.70,4.45);
  \node at (4.375,4.10) {\(j_3\)};
  \draw[dashed,draw=chainOrange!95!black]
    (4.375,4.45)--(4.375,4.85);
  \node[align=center,text=chainOrange!95!black,anchor=north]
    at (4.375,3.68) {same \(r\) as \(i\)\\insert later};
  \draw[<->,draw=chainBlue!90!black,thick]
    (1.25,5.82)--(4.70,5.82);
  \node[anchor=south,text=chainBlue!90!black] at (2.98,5.84)
    {reference distance at most \(D\)};

  % Query-coordinate strip and AVL tree.
  \foreach \x/\q/\fillc in {
      4.95/7/chainGray!10,
      5.75/11/chainBlue!12,
      6.55/14/chainOrange!22,
      7.35/18/chainBlue!12,
      8.15/22/chainBlue!12,
      8.95/27/chainGreen!14}
    {
      \draw[fill=\fillc,draw=black!55]
        (\x,5.08) rectangle (\x+.62,5.62);
      \node at (\x+.31,5.35) {\q};
    }
  \draw[<->,draw=chainOrange!95!black,thick]
    (5.65,4.77)--(8.85,4.77);
  \node[anchor=north] at (7.56,4.68)
    {\(q_i-D<q_j<q_i\)};

  \node[draw=chainBlue!90!black,fill=chainBlue!10,
    minimum width=1.22cm,minimum height=.74cm,align=center]
    (r17) at (7.25,3.70) {\(q=18\)\\max \(h=44\)};
  \node[draw=chainBlue!90!black,fill=chainBlue!10,
    minimum width=1.22cm,minimum height=.74cm,align=center]
    (r11) at (5.75,2.55) {\(q=11\)\\max \(h=44\)};
  \node[draw=chainBlue!90!black,fill=chainBlue!10,
    minimum width=1.22cm,minimum height=.74cm,align=center]
    (r22) at (8.75,2.55) {\(q=22\)\\max \(h=29\)};
  \node[draw=chainOrange!95!black,fill=chainOrange!22,very thick,
    minimum width=1.22cm,minimum height=.70cm,align=center]
    (r14) at (6.35,1.35) {\(q=14\)\\\(h=44\)};
  \node[draw=chainGray!80!black,fill=chainGray!10,
    minimum width=1.22cm,minimum height=.70cm,align=center]
    (r7) at (5.10,1.35) {\(q=7\)\\\(h=20\)};
  \draw[thick] (r17)--(r11);
  \draw[thick] (r17)--(r22);
  \draw[thick] (r11)--(r7);
  \draw[thick] (r11)--(r14);

  % Anchor geometry: the local rectangle and diagonal band are separate tests.
  \node[anchor=west] at (10.15,5.94)
    {local scan if \(j^\star\) is not exact};
  \draw[->] (10.45,3.05)--(14.12,3.05) node[right] {\(r\)};
  \draw[->] (10.45,3.05)--(10.45,5.60) node[left] {\(q\)};
  \begin{scope}
    \clip (10.50,3.10) rectangle (13.90,5.55);
    \fill[chainGreen!12]
      (10.95,3.05)--(11.75,3.05)--(14.25,5.55)--(13.45,5.55)--cycle;
    \draw[dashed,chainGreen!80!black]
      (10.95,3.05)--(13.45,5.55)
      (11.75,3.05)--(14.25,5.55);
  \end{scope}
  \draw[fill=chainBlue!7,fill opacity=.45,draw=chainBlue!90!black,dashed]
    (12.00,3.65) rectangle (13.85,5.50);
  \node[text=chainBlue!90!black] at (12.94,3.88) {local interval};
  \node[anchor=west,text=chainGreen!70!black,align=left]
    at (10.55,3.58) {band\\\(\delta\leq B\)};

  \coordinate (rmqtarget) at (13.85,5.50);
  \coordinate (rmqprimary) at (11.10,4.42);
  \coordinate (rmqlocal) at (12.75,4.40);
  \draw[thick,chainOrange!95!black] (rmqprimary)--(rmqtarget);
  \draw[very thick,chainOrange!95!black]
    (11.92,4.64)--(12.12,4.84) (11.92,4.84)--(12.12,4.64);
  \draw[chain flow,chainGreen!80!black] (rmqlocal)--(rmqtarget);
  \node[circle,draw=chainOrange!95!black,fill=chainOrange!20,
    minimum size=.20cm,inner sep=0pt,label=above:\(j^\star\)]
    at (rmqprimary) {};
  \node[circle,draw=chainBlue!90!black,fill=chainBlue!25,
    minimum size=.20cm,inner sep=0pt,label=left:\(j\)] at (rmqlocal) {};
  \fill[chainBlue!90!black] (13.35,5.02) circle (2pt);
  \draw[very thick,chainOrange!95!black]
    (12.29,5.00)--(12.45,5.16) (12.29,5.16)--(12.45,5.00);
  \node[circle,draw=chainGreen!80!black,fill=chainGreen!20,
    minimum size=.24cm,inner sep=0pt,label=right:\(i\)]
    at (rmqtarget) {};
  \draw[->,thick,chainBlue!90!black] (14.16,5.24)--(14.16,3.82);
  \node[rotate=90,text=chainBlue!90!black] at (14.43,4.52)
    {decreasing \(q_j\)};
  \node[align=center] at (12.30,2.62)
    {\(0<r_i-r_j\leq D_{\mathrm{in}}\)\\
     \(q_i-D_{\mathrm{in}}\leq q_j\leq q_i-1\)};

  \node[draw=chainGreen!80!black,fill=chainGreen!7,
    text width=4.10cm,align=center,inner sep=4pt]
    at (12.30,1.10)
    {For \(j^\star\) and every scanned \(j\):\\
     \(c_j=f_j+E(j,i)\)\\
     update only if \(\delta(j,i)\leq B\) and \(c_j>f_i\)\\
     then \((f_i,p_i)\gets(c_j,j)\)};
\end{tikzpicture}
\caption{Minimap2 RMQ for one target. (a) Retain earlier anchors in the
reference window. (b) Select \(j^\star\); ``max \(h\)'' is the largest
value in a node's subtree. (c) The blue rectangle is the local interval;
the green band satisfies \(\delta\leq B\). Orange crosses mark rejected
candidates, including one inside the local interval. A scanned link
(green arrow) replaces the current result only if both checks below pass.
The local scan can stop early.}
\label{fig:chaining-minimap2-rmq}
\label{fig:chaining-rmq-contrast}
\end{figure}

Algorithm~\ref{alg:chaining-minimap2-rmq} shows the forward pass.
The tree maintenance follows the insertion and removal rules above.
\textsc{RangeBest} returns the anchor with the largest \(h(j)\) in the key
interval, or \(\bot\) if the interval contains no retained anchor.
\textsc{CompleteLink} calculates \(E(j,i)\) from
Equations~\ref{eq:chaining-link-score}--\ref{eq:chaining-penalty}, together
with \(\delta\) and the exact-link condition in
Equation~\ref{eq:chaining-rmq-exact}.
For RMQ scoring, the logarithmic term is zero when \(\delta=0\);
\(\operatorname{mglog}_2(\delta+1)\) is evaluated only when \(\delta\geq1\).

\begin{algorithm}[H]
\caption{Minimap2 RMQ forward chaining}
\label{alg:chaining-minimap2-rmq}
\footnotesize
\begin{algorithmic}[1]
  \Require sorted anchors \(a_0,\ldots,a_{n-1}\), distance \(D\),
    inner distance \(D_{\mathrm{in}}\), bandwidth \(B\), tree cap \(C\),
    skip limit \(S\)
  \Ensure scores \(f\) and predecessor indices \(p\)
  \Function{Minimap2RMQForward}{\(a,D,D_{\mathrm{in}},B,C,S\)}
    \State \(D\gets\max(D,B),\quad
      D_{\mathrm{in}}\gets\min(D,\max(0,D_{\mathrm{in}}))\)
    \State \(T\gets\varnothing,\quad T_{\mathrm{in}}\gets\varnothing\)
          \Comment{augmented AVL trees}
    \State \(t\gets(0,\ldots,0)\)
          \Comment{B1 marks for the optional local scan}
    \For{\(i\gets0\) \textbf{to} \(n-1\)}
      \State \Call{InsertCompletedReferenceGroup}{\(T,a,f,i\)}
      \If{\(D_{\mathrm{in}}>0\)}
        \State \Call{InsertCompletedReferenceGroup}{\(T_{\mathrm{in}},a,f,i\)}
      \EndIf
      \State \Call{RemoveExpiredEntries}{\(T,a,i,D,C\)}
      \If{\(D_{\mathrm{in}}>0\)}
        \State \Call{RemoveExpiredEntries}{\(T_{\mathrm{in}},a,i,D_{\mathrm{in}},C\)}
      \EndIf
      \State \(f_i\gets\ell_i,\quad p_i\gets-1\)
      \State \(j^\star\gets
        \Call{RangeBest}{T,(q_i-D,\mathrm{INT32\_MAX}),(q_i,0)}\)
      \If{\(j^\star\neq\bot\)}
        \State \((e,\delta,\mathrm{exact})\gets
          \Call{CompleteLink}{j^\star,i}\)
              \Comment{\(e=E(j^\star,i)\)}
        \If{\(\delta\leq B\land f_{j^\star}+e>f_i\)}
          \State \(f_i\gets f_{j^\star}+e,\quad p_i\gets j^\star\)
        \EndIf
        \If{\(\neg\mathrm{exact}\land T_{\mathrm{in}}\neq\varnothing
             \land q_i>0\)}
          \State \((f_i,p_i)\gets
            \Call{ScanLocalInterval}{T_{\mathrm{in}},i,D_{\mathrm{in}},B,S,f,p,t}\)
              \Comment{Algorithm~\ref{alg:chaining-rmq-local-scan}}
        \EndIf
      \EndIf
    \EndFor
    \State \Return \(f,p\)
  \EndFunction
\end{algorithmic}
\end{algorithm}

Algorithm~\ref{alg:chaining-rmq-local-scan} expands the local scan.
The mark \(t_j=i\) means that a candidate already tested for target \(i\)
points to \(j\) through its saved predecessor. Using \(i\) as the mark
avoids clearing the array for each new target. A candidate outside the band
does not update the score, counter, or marks. The tree-selected anchor can
be visited again if it also lies in the local interval.

\begin{algorithm}[H]
\caption{Minimap2 local scan after the primary RMQ query}
\label{alg:chaining-rmq-local-scan}
\footnotesize
\begin{algorithmic}[1]
  \Require maintained inner tree \(T_{\mathrm{in}}\), target \(i\),
    local limit \(D_{\mathrm{in}}\), bandwidth \(B\), skip limit \(S\)
  \Require completed earlier scores and predecessors \(f,p\), marks \(t\),
    and current result \((f_i,p_i)\)
  \Ensure updated \((f_i,p_i)\) and marks \(t\)
  \Function{ScanLocalInterval}{\(T_{\mathrm{in}},i,D_{\mathrm{in}},B,S,f,p,t\)}
    \State \(s\gets0\)
    \For{retained \(j\in T_{\mathrm{in}}\) with
      \(q_i-D_{\mathrm{in}}\leq q_j\leq q_i-1\), by decreasing \((q_j,j)\)}
      \State \(c\gets f_j+E(j,i),\quad
        \delta\gets|d_r(j,i)-d_q(j,i)|\)
      \If{\(\delta\leq B\)}
        \If{\(c>f_i\)}
          \State \(f_i\gets c,\quad p_i\gets j,\quad s\gets\max(0,s-1)\)
        \ElsIf{\(t_j=i\)}
          \State \(s\gets s+1\)
          \If{\(s>S\)}
            \State \textbf{break}
          \EndIf
        \EndIf
        \If{\(p_j\geq0\)}
          \State \(t_{p_j}\gets i\)
        \EndIf
      \EndIf
    \EndFor
    \State \Return \((f_i,p_i)\)
  \EndFunction
\end{algorithmic}
\end{algorithm}

Each tree insertion, deletion, and primary query costs
\(O(\log M_{\max})\), where \(M_{\max}\) is the largest number of anchors
retained in a tree. Each anchor is inserted and removed at most once per
tree. If target \(i\) visits \(k_i\) local candidates, the forward cost is
\[
  O(n\log M_{\max})+
  O\!\left(\sum_i k_i\right).
\]
The local scans can still make this quadratic.  The parameter
\texttt{max\_chain\_iter} is not used.  The parameter
\texttt{max\_chain\_skip} affects only the local scan.

Minimap2's manual recommends \texttt{--rmq} for aligning contigs through
long insertions and deletions (INDELs). Its assembly presets enable this
algorithm~\cite{minimap2_rmq_manual}.

\subsection{RMQ re-chaining in the pipeline}
\label{sec:chaining-rmq-rechaining}

\emph{RMQ re-chaining} means running chaining again with a larger bandwidth
to try to connect chains across long gaps. It takes place after the first
backtracking step. For one query, the RMQ re-chaining path follows this order:
\begin{enumerate}
  \item \textbf{Seeding.} Find matching minimizers and form anchors.
  \item \textbf{Initial forward scoring.} Use DP or RMQ to calculate
    scores \(f_i\) and predecessors \(p_i\).
  \item \textbf{Initial backtracking.} Follow those predecessors to build
    chains, applying the minimum score and anchor-count requirements.
    Keep the anchors belonging to the accepted chains.
  \item \textbf{Optional RMQ re-chaining.} Combine these retained anchors
    into one sorted list. Run RMQ forward scoring with a larger bandwidth,
    calculating new scores and predecessors. Then run backtracking again
    to extract new chains. Each input is still an anchor, not a whole
    chain. This step reuses existing anchors and does not repeat seeding.
  \item \textbf{Later mapping stages.} Select and filter candidate mappings,
    then perform base-level alignment if it was requested. Finish mapping
    selection and calculate mapping quality.
\end{enumerate}
This order follows minimap2's chaining and mapping stages
\cite{minimap2_software}.

Minimap2 checks whether to re-chain once after the initial backtracking.
It performs the extra forward scoring and backtracking \textbf{at most once}
for a query against one loaded reference-index part. It does not repeat
these steps until the chains stop changing. With a split reference index, the
query is mapped against each part separately, so each part can make its
own re-chaining decision. There is no setting for additional re-chaining
rounds within one mapping call.

The main command-line controls are listed in
Table~\ref{tab:chaining-rmq-rechaining-options}.
\begin{table}[H]
\centering
\caption{Minimap2 controls for initial chaining and RMQ re-chaining.}
\label{tab:chaining-rmq-rechaining-options}
\small
\setlength{\tabcolsep}{4pt}
\renewcommand{\arraystretch}{1.12}
\begin{tabularx}{\linewidth}{@{}
  >{\raggedright\arraybackslash}p{.28\linewidth}
  >{\raggedright\arraybackslash}X@{}}
\toprule
\textbf{Option} & \textbf{Effect}\\
\midrule
\texttt{-r NUM1,NUM2} & The first value sets bandwidth \(B\) for initial
  chaining. The second sets \(B\) for RMQ re-chaining. Re-chaining requires
  the second value to be larger.\\
\addlinespace
\texttt{--no-long-join} & Disables this RMQ re-chaining step.\\
\addlinespace
\texttt{--rmq=yes} or \texttt{--rmq=no} & Chooses RMQ or DP for initial
  forward scoring. Choosing DP does not disable later RMQ re-chaining.\\
\addlinespace
\texttt{--rmq-inner NUM} & Sets the local distance limit within an RMQ
  forward run. It does not set the number of chaining runs.\\
\bottomrule
\end{tabularx}
\end{table}
For example, the assembly presets use \texttt{-r1k,100k}: bandwidth
\(1000\) for initial chaining and \(100000\) for re-chaining
\cite{minimap2_rmq_manual}. The larger bandwidth permits a larger diagonal
difference. RMQ also raises its distance limit \(D\) to at least \(B\),
so re-chaining can extend the distance window as well.

Re-chaining is conditional even when the options allow it. It requires
more than one retained chain and a single query sequence in the mapping
call. Short-read DNA mapping and spliced alignment skip it.
Minimap2 also checks the query
span of the first retained chain in reference order: the distance between
its first and last anchor on the query. With the default thresholds, either
the query length minus this span must exceed \(1000\) bases, or this span
must exceed \(10\%\) of the query length. These two thresholds can be set
through the library interface, but are not exposed as command-line options
\cite{minimap2_software}.

The local scan belongs inside an RMQ forward run: it compares nearby
predecessors for one target anchor. Re-chaining instead repeats forward
scoring and backtracking over the retained anchor list. BioseqDB's
configured RMQ route performs one forward run and one backtracking step,
with no local scan and no RMQ re-chaining.

\subsection{What can be lost without the local scan?}
\label{sec:chaining-rmq-local-effect}

A larger ranking value can give a smaller complete score. Consider target
\((r_i,q_i)=(100,100)\) and two earlier anchors \(u\) and \(v\), both with
already calculated score \(40\). Let all three seed spans be \(10\),
\(\lambda_g=0.1\), \(\lambda_s=0\), \(D=30\),
\(D_{\mathrm{in}}=20\), and \(B=20\).

\begin{table}[H]
\centering
\caption{Two candidates whose ranking values and complete scores disagree.
The earlier scores are given inputs to this example.}
\label{tab:chaining-rmq-local-example}
\small
\setlength{\tabcolsep}{5pt}
\begin{tabular}{@{}lrrrrrrr@{}}
\toprule
Candidate & \((r_j,q_j)\) & \(f_j\) & \(h(j)\) & \(g\) & \(\delta\)
  & \(E(j,i)\) & \(f_j+E(j,i)\)\\
\midrule
\(u\), tree choice & \((80,95)\) & 40 & 48.75 & 5 & 15 & 2 & 42\\
\(v\), local choice & \((85,85)\) & 40 & 48.50 & 15 & 0 & 10 & 50\\
\bottomrule
\end{tabular}
\end{table}

\begin{itemize}
  \item The primary query chooses \(u\), since \(48.75>48.50\). Its reward
    is \(5\) and its integer gap penalty is \(3\), so it raises
    \(f_i\) from \(10\) to \(42\).
  \item The link from \(u\) is not exact. Both candidates are inside the
    local interval and pass the bandwidth test. If the scan reaches \(v\),
    it raises \(f_i\) to \(50\) and changes \(p_i\) to \(v\).
  \item Without that scan, the result stays at \(42\). Recalculating the
    complete score of \(u\) cannot reveal the better link from \(v\).
\end{itemize}

This is an example of a lost chaining-score improvement, not a measurement
of alignment accuracy. More generally, omitting the local scan can leave a
lower-scoring chain or miss a nearby link when the primary choice is
rejected. Changed scores and predecessors can affect which chains pass
later filters and which regions reach base-level alignment.
The minigraph paper supports using the full gap cost locally
\cite{minimap2_rmq_scoring_minigraph}. The sources reviewed here do not
quantify the alignment-quality loss caused by removing only the local
scan. Its effect on BioseqDB's final alignments remains unmeasured.

\subsection{BioseqDB RMQ differences}
\label{sec:chaining-bioseqdb-rmq-summary}

BioseqDB uses the same tree idea and ranking value \(h(j)\), but its
configured RMQ route sets \(D_{\mathrm{in}}=0\). It therefore calculates
one complete candidate score per target with a nonempty query result.
Although an optional inner scan is implemented, this route does not enable
it.

I gave RMQ lower priority than the ordinary DP chaining described earlier.
In particular, improving chaining across long insertions and deletions was
outside the main implementation work for this thesis. I therefore kept a
simple RMQ version with one primary query per target and no local scan or
RMQ re-chaining (Section~\ref{sec:chaining-rmq-rechaining}).
This is a choice about development scope. It does not establish that
the local scan is unnecessary or that alignment accuracy is unchanged.

\begin{table}[H]
\centering
\caption{Differences between minimap2 and the configured BioseqDB RMQ route.}
\label{tab:chaining-bioseqdb-rmq-differences}
\small
\setlength{\tabcolsep}{4pt}
\renewcommand{\arraystretch}{1.12}
\begin{tabularx}{\linewidth}{@{}
  >{\raggedright\arraybackslash}p{.19\linewidth}
  >{\raggedright\arraybackslash}X
  >{\raggedright\arraybackslash}X@{}}
\toprule
 & \textbf{Minimap2} & \textbf{BioseqDB}\\
\midrule
Local scan & Optional scan with complete scores and the B1 skip rule
  & Disabled by \(D_{\mathrm{in}}=0\)\\
\addlinespace
Query interval & Composite bounds in
  Equation~\ref{eq:chaining-rmq-key-range}
  & Includes both \(q_i-D\) and \(q_i\)\\
\addlinespace
Bandwidth check & Rejects a link when \(\delta>B\)
  & Calculates \(\delta\), but does not reject on this condition\\
\addlinespace
RMQ re-chaining & At most one extra forward run and backtracking step
  with a larger bandwidth & Not performed\\
\addlinespace
Backtracking & Uses the maximum-drop cut search
  & Scalar backtracking without the maximum-drop cut search\\
\bottomrule
\end{tabularx}
\end{table}

Two of these differences are separate correctness limitations:
\begin{itemize}
  \item Including \(q_j=q_i\) allows a link with no query advance. BioseqDB
    searches from \((q_i-D,\mathrm{INT64\_MIN})\) to
    \((q_i,\mathrm{INT64\_MAX})\), so every anchor at either query
    boundary is eligible for selection.
  \item Omitting the bandwidth rejection allows a link farther from the
    diagonal than the configured \(B\). A strictly larger score alone does
    not enforce that limit.
\end{itemize}
Neither behavior follows from disabling the local scan. They must be
considered separately when comparing the two implementations.

BioseqDB also stores coordinates in 64 bits and uses the anchor order in
Equation~\ref{eq:chaining-order}. Minimap2 sorts by its packed reference key.
Together with different rules for choosing between equal ranking values,
this can change which predecessor
is selected when ranking values are equal. With no local scan, neither
\texttt{max\_chain\_skip} nor \texttt{max\_chain\_iter} limits this forward
pass. The configured BioseqDB RMQ route, including sorting and scalar
extraction, takes \(O(n\log n)\) worst-case time.
Section~\ref{sec:chaining-bioseqdb} describes its selection.

\subsection{Future GPU work}

BioseqDB currently has no GPU kernel for RMQ-based chaining. Its Vulkan
chaining algorithm, described next, uses batched DP scoring.

GPU-RMQ by Kreis et al. suggests possible ways to speed up the range-query
part in future work~\cite{gpu_rmq_kreis_2026}:
\begin{itemize}
  \item Build several levels of stored range minima. A query then scans
    only the relevant ranges at each level.
  \item Use groups of GPU threads to read adjacent array entries and find
    their minimum together.
  \item On supported hardware, use dedicated ray-tracing cores for the
    upper level and CUDA cores for the lower levels.
\end{itemize}
The paper studies batches of RMQs on a fixed array. Chaining adds completed
anchors and removes expired ones as scores become available. Applying those
GPU techniques would therefore require a design for these updates and score
dependencies. Faster range queries would also leave the choice of scoring
and local-scan rules to the chaining algorithm. The paper does not establish
a speedup for BioseqDB chaining.
